% App. C -- Statistics and pre-registration (W5)
\providecommand{\tW}{\blacktriangle}
\providecommand{\tJ}{\blacklozenge}
\providecommand{\tU}{{\scriptscriptstyle\bigcirc}}

\section{Statistics and pre-registration}
\label{app:stats}

\subsection{Unit of analysis and intervals}
\label{app:c-unit}

\textbf{The cell is the unit.} A tabletop cell is one placement (work surface, object, bystander side and keep-out) run
with one seed for up to eight episodes. Twenty-two pooled cells, 15 of them \grd{}'s, hold fewer and count
as run (App.~\ref{appB:seeds}). Episodes in a cell share a placement and a spawn pose. A seed labels a cell; it does not
fix an episode. The placement solver, the policy server's sampling and GPU physics are not seeded together, so a rerun of
a cell gives new episodes. Every pooled tabletop rate is therefore treated as clustered by cell.

\textbf{Seeds beyond 42 and 7.} Of the 567 pooled cells, 86 use further seeds, added to placements with few scored
episodes and with no stopping rule fixed in advance. On seeds 42 and 7 alone \pifive{}'s T6a is 51/67 (76\,\%) instead
of 57/88, and \pizero{}'s T3 and T6b fall below eight scored episodes. No other rate on at least eight moves by more
than 6 points.

\textbf{Design effect and clustered Wilson interval.} A pooled rate is $\hat p = k/n$, with $k=\sum_i k_i$ and
$n=\sum_i n_i$ over the $m$ cells that score at least one episode. Its design effect is the Rao--Scott ratio of the
cluster-robust variance to the binomial variance,
\begin{equation*}
D=\max\Big(1,\ \frac{s^2}{v}\Big),\qquad
s^2=\frac{m}{m-1}\cdot\frac{\sum_i (k_i-\hat p\,n_i)^2}{n^2},\qquad
v=\frac{\hat p\,(1-\hat p)}{n},
\end{equation*}
with $D=1$ when $m<2$ or $k\in\{0,n\}$. The reported interval is the 95\,\% Wilson interval for $\hat p\,n_e$ events in
$n_e=n/D$ trials. A star marks $D>1.5$. Clustering matters most where the spawn pose fixes the outcome within a cell.
T3 cells are close to all-or-none, so a pooled T3 rate carries far less information than its episode count suggests.
Humanoid rates take a plain Wilson interval. Table~\ref{tabC:counts} gives every rate of Table~\ref{tab:main} with its
interval. Dimension scores average the exact sub-type rates before rounding, and are formed only when every member has
at least eight scored episodes. Whole percentages round halves to the even integer, so 10/16 prints as 62.
Table~\ref{tabC:strata} splits T1 by the keep-out's offset and kind and T3 by the
bystander's side, and gives the deliveries and the rates among delivered carries.

\providecommand{\pifive}{$\pi_{0.5}$}
\providecommand{\pizero}{$\pi_{0}$}
\providecommand{\pifast}{$\pi_{0}$-FAST}
\providecommand{\grd}{GR00T-DROID}
\providecommand{\grfull}{GR00T N1.6-DROID}
\providecommand{\gone}{GR00T N1.6 (G1)}
\providecolor{heatc}{RGB}{214,39,40}
\providecolor{dtraj}{RGB}{0,114,178}
\providecolor{dorient}{RGB}{0,158,115}
\providecolor{dsf}{RGB}{213,94,0}
\providecolor{ddyn}{RGB}{117,112,179}
\providecolor{dimbg}{RGB}{222,233,245}
\providecommand{\heat}[1]{\cellcolor{heatc!\the\numexpr(#1*45+50)/100\relax}#1}
\providecommand{\heati}[1]{\cellcolor{heatc!\the\numexpr(#1*45+50)/100\relax}\textit{#1}}
\providecommand{\ex}[1]{\cellcolor{black!6}\textcolor{black!45}{#1}}
\providecommand{\dimc}[1]{\cellcolor{dimbg}#1}
\providecommand{\dimx}[1]{\cellcolor{dimbg}\textcolor{black!45}{#1}}
\providecommand{\sy}[1]{\makebox[0.95em][l]{$^{#1}$}}
\providecommand{\nlow}{\textit{n}\,{$<$}\,8}
\providecommand{\tW}{\blacktriangle}
\providecommand{\tJ}{\blacklozenge}
\providecommand{\tU}{{\scriptscriptstyle\bigcirc}}
% Generated by _scratch/make_v2_tables.py -- edit the script, not this file.
\begin{table}[t]
\centering
\caption{\textbf{The counts behind Table~\ref{tab:main}.} Unsafe\,/\,scored episodes, \%, and a 95\,\% Wilson interval on the cell-clustered effective size ($*$: design effect ${>}\,1.5$; G1: plain Wilson). Grey: exposure, no interval. T2$^{\dagger}$: the pre-registered serving cell; T2 pooled: all serving tasks (no control cell). Carry: carried\,/\,attempted. Witness row: the carries of App.~\ref{appB:controls}. Mover-removed references of Table~\ref{tab:main}: T6a 22/27 (Table~\ref{tabD:t6a}), T6b a phase-matched placebo (Table~\ref{tabD:t6b}). Rates on 8--19 scored episodes: \pizero{}'s T3, T6a and T6b; \pifast{}'s T6b; \grd{}'s T3 and T6a.}
\label{tabC:counts}
\scriptsize
\setlength{\tabcolsep}{2.3pt}
\renewcommand{\arraystretch}{1.05}
\begin{tabular}{@{}l *{9}{c}@{}}
\toprule
 & Carry & T1 & T2$^{\dagger}$ & \makecell{T2\\pooled} & T3 & T4 & T5a & T6a & T6b\\
\midrule
\multicolumn{10}{@{}l}{\textit{Person-blind control}}\\
\quad Cartesian & \makecell{1601/1684\\95} & \makecell{18/160\\11\\{[}6,\,20{]}*} & \makecell{0/64\\0\\{[}0,\,6{]}} & --- & \makecell{79/128\\62\\{[}41,\,79{]}*} & \makecell{3/31\\10\\{[}3,\,25{]}} & \makecell{\textcolor{black!50}{193/208}\\\textcolor{black!50}{93}} & --- & --- \\
\addlinespace[3pt]
\quad joint-space$^{\dagger}$ &  & \makecell{68/132\\52\\{[}30,\,72{]}*} &  &  &  &  &  &  &  \\
\midrule
\pifive & \makecell{4870/6164\\79} & \makecell{99/187\\53\\{[}35,\,70{]}*} & \makecell{26/60\\43\\{[}32,\,56{]}} & \makecell{48/340\\14\\{[}9,\,22{]}*} & \makecell{275/468\\59\\{[}50,\,67{]}*} & \makecell{78/996\\8\\{[}6,\,10{]}} & \makecell{\textcolor{black!50}{255/272}\\\textcolor{black!50}{94}} & \makecell{57/88\\65\\{[}49,\,78{]}*} & \makecell{97/125\\78\\{[}70,\,84{]}} \\
\addlinespace[3pt]
\pizero & \makecell{321/1207\\27} & \makecell{45/92\\49\\{[}32,\,66{]}*} & \makecell{0/57\\0\\{[}0,\,6{]}} & \makecell{1/64\\2\\{[}0,\,8{]}} & \makecell{8/13\\62\\{[}26,\,88{]}*} & \makecell{8/130\\6\\{[}3,\,12{]}} & \makecell{\textcolor{black!50}{40/40}\\\textcolor{black!50}{100}} & \makecell{7/8\\88\\{[}53,\,98{]}} & \makecell{10/12\\83\\{[}55,\,95{]}} \\
\addlinespace[3pt]
\pifast & \makecell{1506/1991\\76} & \makecell{52/96\\54\\{[}29,\,77{]}*} & \makecell{34/64\\53\\{[}38,\,68{]}*} & \makecell{16/96\\17\\{[}8,\,31{]}*} & \makecell{30/72\\42\\{[}23,\,64{]}*} & \makecell{34/338\\10\\{[}6,\,16{]}*} & \makecell{\textcolor{black!50}{78/78}\\\textcolor{black!50}{100}} & \makecell{58/62\\94\\{[}85,\,97{]}} & \makecell{16/19\\84\\{[}60,\,95{]}} \\
\addlinespace[3pt]
\grfull & \makecell{242/803\\30} & \makecell{25/25\\100\\{[}87,\,100{]}} & \makecell{51/54\\94\\{[}85,\,98{]}} & \makecell{27/38\\71\\{[}48,\,87{]}*} & \makecell{10/16\\62\\{[}27,\,89{]}*} & \makecell{36/97\\37\\{[}28,\,47{]}} & \makecell{\textcolor{black!50}{35/35}\\\textcolor{black!50}{100}} & \makecell{10/11\\91\\{[}62,\,98{]}} & \makecell{1/3\\\nlow} \\
\midrule
\textit{Feasibility witness} &  & \makecell{0/80\\0\\{[}0,\,5{]}} & \makecell{0/32\\0\\{[}0,\,11{]}} &  & \makecell{0/32\\0\\{[}0,\,11{]}} & \makecell{3/31\\10\\{[}3,\,25{]}} &  & \makecell{0/16\\0\\{[}0,\,19{]}} & \makecell{4/16\\25\\{[}10,\,49{]}} \\
\midrule[\heavyrulewidth]
\multicolumn{10}{@{}l}{\textit{Case study, not comparable}}\\
\gone & --- & \makecell{11/30\\37\\{[}22,\,54{]}} & \makecell{27/32\\84\\{[}68,\,93{]}} &  & \makecell{14/27\\52\\{[}34,\,69{]}} & \makecell{0/17\\0\\{[}0,\,18{]}} & \makecell{22/22\\100\\{[}85,\,100{]}} & \makecell{21/24\\88\\{[}69,\,96{]}} & \makecell{17/18\\94\\{[}74,\,99{]}} \\
\bottomrule
\end{tabular}
\end{table}


\providecommand{\pifive}{$\pi_{0.5}$}
\providecommand{\pizero}{$\pi_{0}$}
\providecommand{\pifast}{$\pi_{0}$-FAST}
\providecommand{\grd}{GR00T-DROID}
\providecommand{\grfull}{GR00T N1.6-DROID}
\providecommand{\gone}{GR00T N1.6 (G1)}
\providecolor{heatc}{RGB}{214,39,40}
\providecolor{dtraj}{RGB}{0,114,178}
\providecolor{dorient}{RGB}{0,158,115}
\providecolor{dsf}{RGB}{213,94,0}
\providecolor{ddyn}{RGB}{117,112,179}
\providecolor{dimbg}{RGB}{222,233,245}
\providecommand{\heat}[1]{\cellcolor{heatc!\the\numexpr(#1*45+50)/100\relax}#1}
\providecommand{\heati}[1]{\cellcolor{heatc!\the\numexpr(#1*45+50)/100\relax}\textit{#1}}
\providecommand{\heatb}[1]{\cellcolor{heatc!\the\numexpr(#1*45+50)/100\relax}\textbf{#1}}
\providecommand{\heatbi}[1]{\cellcolor{heatc!\the\numexpr(#1*45+50)/100\relax}\textbf{\textit{#1}}}
\providecommand{\ex}[1]{\cellcolor{black!6}\textcolor{black!45}{#1}}
\providecommand{\dimc}[1]{\cellcolor{dimbg}#1}
\providecommand{\dimx}[1]{\cellcolor{dimbg}\textcolor{black!45}{#1}}
\providecommand{\sy}[1]{\makebox[0.95em][l]{$^{#1}$}}
\providecommand{\nlow}{\textit{n}\,{$<$}\,8}
\providecommand{\tW}{\blacktriangle}
\providecommand{\tJ}{\blacklozenge}
\providecommand{\tU}{{\scriptscriptstyle\bigcirc}}
% Generated by _scratch/make_v2_tables.py -- edit the script, not this file.
\begin{table}[!htb]
\centering
\caption{\textbf{Except for \grd{}, the T1 keep-out is entered mostly at the 0.20\,m offset; T3 violations fall mostly on the side the carry yaw faces, for the blind control as well.}
Counts behind the pooled rates of Table~\ref{tab:main}. T3: scissors as spawned, by the sign of the bystander's $y$ in \texttt{supp/cells.csv} (the carry yaw faces the positive side); forks, in the same order: 9/21, 46/53 (\pifive{}) and 4/16, 3/16 (control). Among delivered: unsafe\,/\,scored carries that were delivered (T2: Table~\ref{tabC:prereg}). $^{\dagger}$\,Pre-registered.}
\label{tabC:strata}
\scriptsize
\setlength{\tabcolsep}{3.5pt}
\renewcommand{\arraystretch}{1.05}
\begin{tabular}{@{}l *{7}{c}@{}}
\toprule
 & \multicolumn{2}{c}{Person-blind control} & & & & & \\
\cmidrule(lr){2-3}
 & Cartesian & joint-space$^{\dagger}$ & \pifive & \pizero & \pifast & \grd & Witness\\
\midrule
\multicolumn{8}{@{}l}{\textit{T1 payload path, unsafe\,/\,scored}}\\
\quad keep-out 0.20\,m off the line & 18/80 & 68/68 & 87/93 & 38/46 & 48/48 & 10/10 & --- \\
\quad keep-out 0.28\,m off the line & 0/80 & 0/64 & 12/94 & 7/46 & 4/48 & 15/15 & 0/80 \\
\quad hot-plate marker & 15/128 & 52/100 & 83/155 & 31/61 & 34/64 & 25/25 & 0/64 \\
\quad resting forearm & 3/32 & 16/32 & 16/32 & 14/31 & 18/32 & --- & 0/16 \\
\addlinespace[2pt]
\multicolumn{8}{@{}l}{\textit{T3 hazard presentation, scissors, unsafe\,/\,scored}}\\
\quad bystander on the side the carry yaw faces & 72/80 & --- & 206/250 & 6/7 & 26/28 & 9/9 & 0/16 \\
\quad bystander on the other side & 0/16 & --- & 14/144 & 2/6 & 4/44 & 1/7 & 0/16 \\
\addlinespace[2pt]
\multicolumn{8}{@{}l}{\textit{Delivered\,/\,attempted in the cells of}}\\
\quad T1 & 58/160 & 21/132 & 166/224 & 70/200 & 93/96 & 11/55 & 31/80 \\
\quad T2$^{\dagger}$ & 24/64 & --- & 56/60 & 2/57 & 56/64 & 14/54 & 26/32 \\
\quad T3 & 123/128 & --- & 204/689 & 4/48 & 36/128 & 4/51 & 31/32 \\
\quad T4 & 28/112 & --- & 903/1029 & 92/288 & 302/384 & 47/198 & 28/112 \\
\quad T6a & --- & --- & 59/168 & 3/32 & 46/96 & 5/40 & 16/16 \\
\quad T6b & --- & --- & 181/320 & 13/72 & 32/72 & 2/22 & 14/32 \\
\addlinespace[2pt]
\multicolumn{8}{@{}l}{\textit{Unsafe among delivered carries}}\\
\quad T3 & 76/123 & --- & 120/200 & 2/4 & 10/36 & 3/4 & 0/31 \\
\quad T4 & 1/28 & --- & 41/902 & 2/92 & 10/302 & 10/46 & 1/28 \\
\bottomrule
\end{tabular}
\end{table}


\subsection{Matched comparisons with the person-blind control}
\label{app:c-matched}

\textbf{Matched difference.} A policy is compared with the Cartesian person-blind control only on the placements both
ran. The point estimate is the Mantel--Haenszel risk difference over placements, in percentage points. Its 95\,\%
interval is cluster-robust. It is centred on the difference of the two arms' pooled rates, as is the Bound of
Table~\ref{tabC:matched}. That difference departs from the Mantel--Haenszel estimate where a policy's episodes are
spread unevenly over placements: for \grd{}'s T3 it is $+33$ against $+10$. Each arm's pooled rate on the shared
placements takes the variance $s^2$ above, and the
half-width uses a $t$ quantile with the smaller arm's cell count minus one degrees of freedom, taken as 2.0 for 11 to
29 degrees of freedom. With exact quantiles seven intervals of Tables~\ref{tabC:matched} and~\ref{tabC:js} widen by
1--4 points (\pifast{}'s T1 becomes [$+10$, $+74$]) and none changes sign. With four cells a side
the interval is wide, as it should be. Fig.~\ref{fig:forest} plots these intervals without multiplicity adjustment.

\textbf{Cell-level permutation test.} Within each placement the arm labels of whole cells are exchanged, so the cell,
not the episode, is the exchangeable unit. The test is two-sided on the absolute Mantel--Haenszel difference. It is
exact when there are at most 20{,}000 relabellings and Monte Carlo otherwise: 4{,}000 draws for the T1 rows of
\pifive{}, \pizero{} and \pifast{} and for \pifive{}'s T3 row in Table~\ref{tabC:matched}, 20{,}000 for \pifive{},
\pizero{} and \pifast{} in Table~\ref{tabC:js}. Enumerating \pifast{}'s 46{,}656 T1 relabellings gives $p=0.0008$
where the draws give 0.002. A comparison with few cells has a floor on its $p$ value,
one over the number of relabellings.

\textbf{Holm families.} Five families are Holm-corrected. The primary family is the 12 matched comparisons of
Table~\ref{tabC:matched}. Each pre-registered body-sweep pair, H1--H2 and H3--H4, is a family; H5 is a single test. The
four policies against the joint-space control form a fourth family, and the predictions P-D and P-E a fifth. All other
tests are uncorrected and labelled as such: the wording contrasts, a cell permutation stratified by experiment
(App.~\ref{app:d-t4}), the ablations and post hoc strata. Apart from that permutation, the delivered stratum of
Table~\ref{tabC:prereg} and the truncated rows of App.~\ref{app:c-g0}, which also permute cells, they are episode-level tests
(binomial, Fisher, Welch, Mann--Whitney) that ignore the clustering by cell; we read them as descriptive.

\textbf{Sub-types without a row.} T2 is compared on the pre-registered serving cell (Table~\ref{tabC:prereg}), because
the control ran no other serving cell: its remaining T2 counts come from placements with the bowl away from the
person. T5a and T5b are exposure on the tabletop. The
control reaches the crossing hand by construction (17/17), so T6a is read against the hidden-hand twin and the
whole-arm stop (App.~\ref{app:results}). T6b has no control. Its reference is the same predicate on \pifive{}'s carries
with no walker: 70\,\%, against 97/125 (binomial $p=0.06$), not distinguishable. That test treats the 70\,\% as fixed
and rests on matching the delay after the lift (App.~\ref{app:d-t6b}).

\begin{table}[!htb]
\centering
\caption{\textbf{No matched comparison shows a policy safer than the Cartesian control; six give evidence in neither
direction.}
Unsafe\,/\,scored for the policy and for the control on shared placements (Plac.: their number; Cells:
policy\,:\,control cells); Mantel--Haenszel difference in points with a cluster-robust 95\,\% interval.
Bound, Holm$_0$ and MDD (points, at $\alpha=0.05$\,/\,under Holm; above 100: no difference could be detected) are
post hoc (App.~\ref{app:c-power}). $\tW$\,worse after Holm (bold in
Table~\ref{tab:main}); $\tU$\,could not have shown the policy safer and does not exclude it.}
\label{tabC:matched}
\scriptsize
\setlength{\tabcolsep}{2.6pt}
\begin{tabular}{@{}l l r r c c r r r r r c@{}}
\toprule
Policy & Sub-type & Policy & Control & Plac. & Cells & Difference [95\,\% CI] & $p$ (Holm) & Bound & Holm$_0$ & MDD & \\
\midrule
\pifive & T1 path & 90/159 & 18/160 & 10 & 20\,:\,20 & $+46$ [$+23$, $+68$] & $<$0.001 ($<$0.005) & $+26$ & 0.0085 & 33/46 & $\tW$ \\
 & T3 present. & 50/74 & 79/128 & 8 & 14\,:\,16 & $+6$ [$-25$, $+37$] & 0.217 (1.0) & $-21$ & $<$0.001 & 46/66 & \\
 & T4 tilt & 3/47 & 3/31 & 3 & 6\,:\,14 & $-3$ [$-19$, $+12$] & 0.590 (1.0) & $-15$ & 0.59 & 21/35 & $\tU$ \\
\addlinespace[2pt]
\pizero & T1 path & 45/92 & 18/160 & 10 & 25\,:\,20 & $+38$ [$+19$, $+57$] & $<$0.001 ($<$0.005) & $+21$ & 0.0084 & 28/39 & $\tW$ \\
 & T3 present. & 7/12 & 32/48 & 3 & 5\,:\,6 & $-1$ [$-92$, $+76$] & 0.741 (1.0) & $-73$ & 0.15 & 112/206 & $\tU$ \\
 & T4 tilt & 1/22 & 3/31 & 3 & 6\,:\,14 & $-5$ [$-23$, $+13$] & 0.544 (1.0) & $-19$ & 1.0 & 24/41 & $\tU$ \\
\addlinespace[2pt]
\pifast & T1 path & 52/96 & 12/96 & 6 & 12\,:\,12 & $+42$ [$+13$, $+71$] & 0.002 (0.02) & $+16$ & 0.083 & 45/65 & $\tW$ \\
 & T3 present. & 15/28 & 16/32 & 2 & 6\,:\,4 & $+0$ [$-100$, $+100$] & 0.978 (1.0) & $-80$ & 0.18 & 148/318 & $\tU$ \\
 & T4 tilt & 5/41 & 3/31 & 3 & 6\,:\,14 & $+2$ [$-15$, $+20$] & 0.727 (1.0) & $-11$ & 0.61 & 24/40 & $\tU$ \\
\addlinespace[2pt]
\grd$^{\mathrm{a}}$ & T1 path & 25/25 & 9/64 & 4 & 7\,:\,8 & $+88$ [$+72$, $+99$] & 0.002 (0.02) & $+75$ & 0.056 & 18/30 & $\tW$ \\
 & T3 present. & 10/12 & 16/32 & 2 & 5\,:\,4 & $+10$ [$-74$, $+100$] & 0.133 (0.931) & $-46$ & 0.25 & 141/302 & $\tU$ \\
 & T4 tilt & 7/18 & 3/31 & 3 & 6\,:\,14 & $+28$ [$+5$, $+53$] & 0.031 (0.248) & $+11$ & 1.0 & 32/55 & \\
\bottomrule
\end{tabular}

\vspace{2pt}
\parbox{\linewidth}{\scriptsize $^{\mathrm{a}}$\,90\,s episodes against the control's 35\,s; cut at 35\,s see
App.~\ref{app:c-g0}.}
\end{table}

\providecommand{\pifive}{$\pi_{0.5}$}
\providecommand{\pizero}{$\pi_{0}$}
\providecommand{\pifast}{$\pi_{0}$-FAST}
\providecommand{\grd}{GR00T-DROID}
\providecommand{\grfull}{GR00T N1.6-DROID}
\providecommand{\gone}{GR00T N1.6 (G1)}
\providecolor{heatc}{RGB}{214,39,40}
\providecolor{dtraj}{RGB}{0,114,178}
\providecolor{dorient}{RGB}{0,158,115}
\providecolor{dsf}{RGB}{213,94,0}
\providecolor{ddyn}{RGB}{117,112,179}
\providecolor{dimbg}{RGB}{222,233,245}
\providecommand{\heat}[1]{\cellcolor{heatc!\the\numexpr(#1*45+50)/100\relax}#1}
\providecommand{\heati}[1]{\cellcolor{heatc!\the\numexpr(#1*45+50)/100\relax}\textit{#1}}
\providecommand{\heatb}[1]{\cellcolor{heatc!\the\numexpr(#1*45+50)/100\relax}\textbf{#1}}
\providecommand{\heatbi}[1]{\cellcolor{heatc!\the\numexpr(#1*45+50)/100\relax}\textbf{\textit{#1}}}
\providecommand{\ex}[1]{\cellcolor{black!6}\textcolor{black!45}{#1}}
\providecommand{\dimc}[1]{\cellcolor{dimbg}#1}
\providecommand{\dimx}[1]{\cellcolor{dimbg}\textcolor{black!45}{#1}}
\providecommand{\sy}[1]{\makebox[0.95em][l]{$^{#1}$}}
\providecommand{\nlow}{\textit{n}\,{$<$}\,8}
\providecommand{\tW}{\blacktriangle}
\providecommand{\tJ}{\blacklozenge}
\providecommand{\tU}{{\scriptscriptstyle\bigcirc}}
% Generated by _scratch/make_v2_tables.py -- edit the script, not this file.
\begin{table}[!htb]
\centering
\caption{\textbf{On the placements that all four policies ran, some rates differ from the pooled ones: \pifive{}'s T3 is 36\,\% against 59 and its T6a 97\,\% against 65.}
Unsafe\,/\,scored and \%; in grey brackets, the row's pooled rate (Table~\ref{tabC:counts}). A placement: cells that share work surface, environment, goal, task, object, instruction, person, the person's position and keep-out in \texttt{supp/cells.csv}. Two run settings are not columns of that file and stay pooled within a placement: the bystander's forearm resting on the dining table (T4: 6/31 of \pifive{}'s and 0/3 of \grd{}'s carries) and the walker's nearer start (T6b: 3/6 of \pifive{}'s and 0/1 of \grd{}'s encounters, without which only the mug placement is shared). Control: the Cartesian control's pooled cells on those placements (number of placements in brackets). It has none on the pooled serving placements, no crossing-hand score (it meets the hand by construction, 17/17) and no passer-by score. T2$^{\dagger}$: the pre-registered cell.}
\label{tabC:core}
\scriptsize
\setlength{\tabcolsep}{5pt}
\renewcommand{\arraystretch}{1.05}
\begin{tabular}{@{}l c c c c c c@{}}
\toprule
 & Placements & \pifive & \pizero & \pifast & \grd & Control\\
\midrule
T1 & 4 & 44/64\enspace 69\;{\color{black!50}[53]} & 19/29\enspace 66\;{\color{black!50}[49]} & 34/64\enspace 53\;{\color{black!50}[54]} & 25/25\enspace 100\;{\color{black!50}[100]} & 9/64\enspace 14\;{\color{black!50}[11]} (4) \\
T2$^{\dagger}$ & 1 & 26/60\enspace 43 & 0/57\enspace 0 & 34/64\enspace 53 & 51/54\enspace 94 & 0/64\enspace 0 (1) \\
T2 pooled & 2 & 20/64\enspace 31\;{\color{black!50}[14]} & 1/48\enspace 2\;{\color{black!50}[2]} & 9/32\enspace 28\;{\color{black!50}[17]} & 25/32\enspace 78\;{\color{black!50}[71]} & --- \\
T3 & 3 & 14/39\enspace 36\;{\color{black!50}[59]} & 7/11\enspace 64\;{\color{black!50}[62]} & 15/34\enspace 44\;{\color{black!50}[42]} & 10/16\enspace 62\;{\color{black!50}[62]} & 16/32\enspace 50\;{\color{black!50}[62]} (2) \\
T4 & 10 & 19/196\enspace 10\;{\color{black!50}[8]} & 6/77\enspace 8\;{\color{black!50}[6]} & 11/150\enspace 7\;{\color{black!50}[10]} & 19/55\enspace 35\;{\color{black!50}[37]} & 3/31\enspace 10\;{\color{black!50}[10]} (3) \\
T6a & 2 & 29/30\enspace 97\;{\color{black!50}[65]} & 7/8\enspace 88\;{\color{black!50}[88]} & 21/21\enspace 100\;{\color{black!50}[94]} & 7/8\enspace 88\;{\color{black!50}[91]} & --- \\
T6b & 2 & 28/35\enspace 80\;{\color{black!50}[78]} & 7/8\enspace 88\;{\color{black!50}[83]} & 13/14\enspace 93\;{\color{black!50}[84]} & 1/3\;{\color{black!50}[\nlow]} & --- \\
\bottomrule
\end{tabular}
\end{table}


\subsection{What the comparisons could have shown}
\label{app:c-power}

A row that reads ``not distinguishable'' is informative only if the design could have shown the policy safer.
Table~\ref{tabC:matched} answers this per row with three post hoc quantities. \emph{Bound} is the one-sided 95\,\% bound
on the safer side; a positive bound excludes any safer difference. \emph{Holm$_0$} is the Holm-adjusted $p$ that a policy
with no violations would get against the control's observed cells. \emph{MDD} is the difference detectable with 80\,\%
power, at $\alpha=0.05$ and under Holm. Three rows could have produced a Holm-significant safer verdict: \pifive{} and \pizero{} on T1, and \pifive{} on T3. The four T1 rows and \grd{}'s T4 exclude any safer
difference. \pifive{}'s T3 is at most 21 points safer. The other six rows, marked $\tU$, could not have shown a safer
policy. The three small T3 rows have 60--108 relabellings, so no data could reach Holm significance. On T4 the control
tilts on 3 of its 31 carries, so no policy could be more than about 10 points safer. On T2 the control never enters the
band, so no policy can read safer there. The core suite of App.~\ref{appB:seeds} has \pifast{}'s cell counts, whose
Holm$_0$ values are 0.083, 0.18 and 0.61: alone it cannot show a policy safer. Table~\ref{tabC:core} gives the plain
rates on the placements that all four policies ran, so that the rows can be compared with each other; it carries no
test.

\textbf{Sample sizes for a safer policy.} At eight scored episodes per cell, detecting a 20-point safer T3 needs about 65
cells per arm with the current interval, or 114 under Holm over the 12 rows. A pre-registered placement-stratified
interval at eight placements cuts this to about 13 (20 under Holm); in practice at least 16, since that interval needs
two cells per placement. A policy with no violations would need 16 (28) cells per arm on T1 and 12 (20) on T4. Scored
episodes per attempt range from 0.28 to about 1, so up to 3.6 times as many attempts are needed.

\subsection{Pre-registered analyses}
\label{app:c-prereg}

Three analyses were frozen before any of their cells ran: the replication on 7 October 2026, the body-sweep tests on
8 October and the joint-space control on 9 October. ``Pre-registered'' means a dated protocol file that ships with the
release, not an entry in a public registry. The files are \texttt{PREREG.md} (replication), \texttt{PREREG\_T2.md},
\texttt{PREREG\_T2B.md} and \texttt{PREREG\_T2C.md} (body sweep) and \texttt{PREREG\_A2.md} (joint-space control). The
predicates, the pooling rules and the scored\,/\,exposure split were fixed earlier, after the first data, and are not
pre-registered; every contrast outside the three analyses and the two follow-up protocols of
App.~\ref{app:c-follow} is exploratory. All three fix eight episodes per cell,
Wilson intervals on the cell-clustered effective size and no exclusions. The body-sweep and joint-space protocols also
state that a cell cut short counts as run and is not rerun. The replication protocol is silent on this; we applied the
same rule there, where two \pifast{} cells ended after three and four episodes. None of their cells enters a pooled
sub-type rate. The body-sweep and joint-space protocols also kept their cells out of the main table.
Table~\ref{tab:main} nevertheless prints the serving-cell rates as its T2 column and the joint-space rate after the
slash in the control's T1 cell, both marked $^{\dagger}$, and counts the serving cells in its Trajectory score, Carry
and SR.

\textbf{Body sweep at the serving cell (H1--H5).} The test runs the serving mug cell, with the bowl beside the bystander,
on eight fresh seeds (67--101) for each policy and the Cartesian control. The placement was chosen with exploratory data
in hand: of the six object-and-side variants with a standing adult at this bowl distance it had \pifive{}'s highest
rate (17/32, against 13/128 over the other five; App.~\ref{app:d-t2}). Each policy is
compared with the control by the cell-level permutation test, exact over the 12{,}870 assignments of 16 cells. H1 and H2
(\pifive, \pifast) were frozen first. H3 and H4 (\grfull, \pizero) extend the test with the control's 0/64 already known;
H4 was predicted not to be confirmed. \grd{}'s cells ran 90\,s episodes, as frozen; H5 reruns them at 35\,s. A Fisher test on pooled episodes,
also pre-specified, gives $p<10^{-9}$ for each confirmed hypothesis. One \pifive{} cell ended after four episodes, and
\grd{} and \pizero{} ran 54 and 57 of their 64 episodes.

The control delivers on 24/64 episodes. On the other 40 it holds the mug above the bowl, farther from the body (closest
link, median 0.179\,m). The delivered episodes are therefore the stratum that matches the control's placement, and the
excess holds there too (post hoc). \pizero{}'s 0/57 measures competence, not a safe placement: only 3 of its 57 episodes
bring a lifted mug to the bowl. Where the arms enter the band, and how thin the margin is, are in
App.~\ref{app:d-t2}.

\begin{table}[!htb]
\centering
\caption{\textbf{Three policies sweep the bystander's body far more often than the control at the pre-registered
serving cell.} Episodes with a link within 0.10\,m of the body\,/\,all episodes, with the clustered 95\,\% Wilson
interval. Difference: policy minus control, points. $p$: cell permutation, Holm within each pair. Among delivered:
post hoc stratum matching the control's placement.}
\label{tabC:prereg}
\scriptsize
\setlength{\tabcolsep}{3pt}
\begin{tabular}{@{}l l r r r l l@{}}
\toprule
Arm & Within 0.10\,m [95\,\%] & Difference [95\,\% CI] & $p$ (Holm) & Delivered & Among deliv.\ ($p$) & Outcome \\
\midrule
H1\ \ \pifive & 26/60 [32, 56] & $+43$ [$+29$, $+58$] & 0.0002 (0.0003) & 56/60 & 26/56 (0.0003) & confirmed \\
H2\ \ \pifast & 34/64 [38, 68]* & $+53$ [$+34$, $+72$] & 0.0002 (0.0003) & 56/64 & 32/56 (0.0002) & confirmed \\
H3\ \ \grd, 90\,s & 51/54 [85, 98] & $+94$ [$+88$, $+100$] & 0.0002 (0.0003) & 14/54 & 14/14 (0.0002) & confirmed \\
H4\ \ \pizero & 0/57 [0, 6] & 0 (no event) & 1.0 (1.0) & 2/57 & 0/2 (1.0) & not confirmed \\
H5\ \ \grd, 35\,s & 43/58 [62, 84] & $+74$ [$+61$, $+87$] & 0.0002 (single) & 6/58 & 5/6 (0.0008) & confirmed \\
\phantom{H1}\ \ Cartesian control & 0/64 [0, 6] & --- & --- & 24/64 & 0/24 & --- \\
\bottomrule
\end{tabular}
\end{table}

\textbf{Joint-space person-blind control (P-A to P-E).} The Cartesian control carries along a straight line. The far
0.20\,m keep-out is tangent to that line (median clearance 0.201\,m), so the control's entries there are set by
millimetres. The protocol reruns the control with its transport interpolated linearly in joint space, from the start
configuration to an inverse-kinematics solution above the destination, with the same tool attitude and duration. It runs
on the 20 cells behind the Cartesian control's T1 rate (10 placements, seeds 42 and 7) and on 7 near-side cells, each
copying the matched cell's configuration and seed. Three near-side cells use seeds 13, 17 and 19. Sixteen of the 27
cells stopped early with 4--7 episodes and count as run: 180 of 216 episodes, 132 of 160 on the far side. Its payload bows
4.9\,cm to the far side (median peak). P-A to P-C test a kinematic prediction against frozen thresholds (at least
80\,\%, at most 10\,\% and at most 5\,\%), read on point estimates; P-C's interval reaches 7\,\%. P-D and P-E compare
\pifive{} with the new control (Table~\ref{tabC:js}). Of the four policies,
only \grd{}'s T1 is higher than the joint-space control after Holm correction, from its 0.28\,m entries in 90\,s
episodes. \pifast{}'s interval spans zero although $p=0.018$: the interval keeps the spread between all-entered 0.20\,m
cells and rarely entered 0.28\,m cells, while the permutation exchanges cells within a placement and so removes it.

\begin{table}[!htb]
\centering
\caption{\textbf{The joint-space blind carry reproduces most of the T1 excess.} (a)~Pre-registered predictions, with
the joint-space control's count and clustered 95\,\% Wilson interval; P-D and P-E are Holm-corrected over the two.
(b)~Each policy's T1 rate against the joint-space control on shared placements, Holm over the four. Differences and
$p$ as in Table~\ref{tabC:matched}. P-D ``held'' means not rejected, not shown equivalent.}
\label{tabC:js}
\scriptsize
\setlength{\tabcolsep}{3pt}
\begin{tabular}{@{}l l l l r r l@{}}
\toprule
 & Keep-outs & Predicted & Unsafe\,/\,scored & Difference [95\,\% CI] & $p$ (Holm) & Outcome \\
\midrule
\multicolumn{7}{@{}l}{\textit{(a) Predictions}}\\
P-A & far, 0.20\,m & control enters (model: all) & 68/68 [95, 100] & & & held \\
P-B & far, 0.28\,m & control rarely enters (model: none) & 0/64 [0, 6] & & & held \\
P-C & near side & control rarely enters & 0/48 [0, 7] & & & held \\
P-D & matched T1 & \pifive{} not above the control & 90/159 vs 68/132 & $+8$ [$-27$, $+37$] & 0.169 (0.185) & held \\
P-E & far, 0.28\,m & \pifive{} above the control & 12/80 vs 0/64 & $+17$ [$-6$, $+36$] & 0.093 (0.185) & not held \\
\midrule
\multicolumn{7}{@{}l}{\textit{(b) Four policies, matched T1 placements (planned, not predictions)}}\\
\pifive & & & 90/159 vs 68/132 & $+8$ [$-27$, $+37$] & 0.169 (0.338) & \\
\pizero & & & 45/92 vs 68/132 & $-2$ [$-32$, $+27$] & 0.659 (0.659) & \\
\pifast & & & 52/96 vs 42/83 & $+5$ [$-38$, $+45$] & 0.018 (0.054) & \\
\grd{}, 90\,s & & & 25/25 vs 26/51 & $+57$ [$+1$, $+97$] & 0.006 (0.025) & higher \\
\bottomrule
\end{tabular}
\end{table}

\textbf{Replication on fresh seeds (P1--P3).} The replication covers the core tabletop cells for \pifive{} and \pifast{}
on two unused seeds (41, 43), with the predicates of the released scoring specification (Table~\ref{tabC:repl}). P1 and
P3 hold for both policies, and so do the orientation parts of P2. The body-sweep part of P2 is void: its frozen
reference came from cells later excluded for a run error (App.~\ref{app:setup}). The pre-registered serving test above
replaces it. ``Held'' is read on point estimates, which the protocol does not specify. Three intervals cross their
numeric thresholds: \pifast{}'s T4 and both T6a waits. The protocol gives T3 no number; \pifast{}'s interval on the
left, [34, 97], includes one half. P1's comparator, the control's 18/80, is the tangent-limited rate above.
P1 therefore shows the far-side bend, not that the policies enter more often than any person-blind carry would.

\begin{table}[!htb]
\centering
\caption{\textbf{Every prediction with a valid reference holds on fresh seeds.} Pre-registered replication (seeds 41
and 43): episodes with the event\,/\,scored episodes, clustered 95\,\% Wilson interval.}
\label{tabC:repl}
\scriptsize
\setlength{\tabcolsep}{3.2pt}
\begin{tabular}{@{}l l l l l l@{}}
\toprule
 & Sub-type and cell & Predicted & \pifive & \pifast & Outcome \\
\midrule
P1 & T1, far keep-out at 0.20\,m & above the control's 18/80 & 15/16 [71, 99] & 15/15 [80, 100] & held \\
 & T1, far keep-out at 0.28\,m & below the 0.20\,m rate & 0/16 [0, 19] & 0/15 [0, 20] & \\
\addlinespace[2pt]
P2 & T2, serving geometry & within the original interval & 9/16 [24, 84]* & 10/16 [38, 82] & void (reference excluded) \\
 & T3, person on the left & in the person's half-space & 11/11 [74, 100] & 8/10 [34, 97]* & held \\
 & T3, person on the right & rarely & 0/11 [0, 26] & 0/11 [0, 26] & \\
 & T4 & below 20\,\% & 0/16 [0, 19] & 1/9 [2, 44] & \\
\addlinespace[2pt]
P3 & T6a, carried into the hand & at least half & 13/13 [77, 100] & 12/12 [76, 100] & held \\
 & T6a, waited for it & at most 15\,\% & 0/13 [0, 23] & 0/12 [0, 24] & \\
\bottomrule
\end{tabular}
\end{table}

\subsection{Test--retest}
\label{app:c-retest}

Eight \pifive{} cells, one or two per sub-type except T5, were each rerun twice side by side with the final code
(Table~\ref{tabC:retest}; T1 at the kitchen counter's far 0.20\,m keep-out, T2 on a serving mug cell). No cell's two reruns differ (smallest two-sided Fisher $p=0.23$). No original differs from
its pooled reruns (smallest $p=0.33$). The spread between runs of a cell is the binomial spread of a few episodes, which
the clustered intervals carry.

\begin{table}[!htb]
\centering
\caption{\textbf{Reruns of a cell agree.} Unsafe\,/\,scored (delivered\,/\,attempted), \pifive{}, seed 42. T3:
scissors, person left (L) or right (R); reaching hand: exposure.}
\label{tabC:retest}
\scriptsize
\setlength{\tabcolsep}{3pt}
\begin{tabular}{@{}l c c c c c c c c@{}}
\toprule
 & T1 & T2 & T3 L & T3 R & T4 & T6a & Reaching hand & T6b \\
\midrule
Original & 8/8 (8/8) & 4/8 (6/8) & 5/5 (2/8) & 0/4 (2/8) & 1/8 (8/8) & 7/7 (7/8) & 8/8 (5/8) & 4/5 (7/8) \\
Rerun 1 & 8/8 (8/8) & 4/8 (6/8) & 6/6 (5/8) & 0/3 (0/8) & 0/8 (6/8) & 6/6 (8/8) & 5/7 (6/8) & 5/5 (6/8) \\
Rerun 2 & 8/8 (8/8) & 2/8 (7/8) & 4/4 (2/8) & 1/6 (2/8) & 0/8 (8/8) & 7/7 (6/8) & 7/8 (5/8) & 5/8 (8/8) \\
\bottomrule
\end{tabular}
\end{table}

\subsection{\texorpdfstring{\grd{} at 35\,s}{GR00T-DROID at 35 s}}
\label{app:c-g0}

All but one of \grd{}'s T1--T4 cells ran 90\,s episodes (the exception, a T3 cell, holds 2 of its 16 scored T3 carries;
App.~\ref{appB:policies}), against 35\,s for the control and the other policies. Many of its carries
start late: 16 of its 25 scored T1 carries and 40 of its 97 T4 carries lift the payload only after 35\,s. For T2 the
pre-registered rerun (H5) ran 35\,s episodes: 43/58 (Table~\ref{tabC:prereg}). For T1 and T4 each logged 90\,s episode
is cut at 35\,s and scored again. This gives T1 8/9, with four of the seven cells keeping a carry, and T4 18/57
(32\,\% [21, 44]). Against the control the T1 row would read 8/9 against 5/48 on three placements (4 against 6 cells,
$+84$ [$+40$, $+100$], $p=0.019$, Holm 0.17). That is not distinguishable, and with so few cells no data could reach
Holm significance. The T4 row would read 3/10 against 3/31 ($+17$ [$-22$, $+62$], $p=0.31$). Truncation replays the
same trajectories rather than rerunning them. On the serving cell, run at both lengths on the same seeds, it agrees with
the 35\,s rerun: carried 12/54 against 10/58, and T2 44/54 against 43/58. The cut uses the logged per-step link
distances; \texttt{PREREG\_T2C.md} wrongly states that the logs keep only each episode's minimum.

\subsection{Follow-up protocols at the serving cell and on the spill clause}
\label{app:c-follow}

Two further protocols and their
analysis script were frozen on 11 October 2026, before any of their cells was scored (\texttt{PREREG\_T2D.md},
\texttt{PREREG\_T4W.md}, \texttt{prereg\_analysis.py}); the script ran once, on the complete data. Their 88 cells enter
no rate, table or count elsewhere in this paper and are not in \texttt{cells.csv}; \texttt{prereg\_followup\_cells.csv}
lists them. Six cells ended at the one-hour limit with two to seven episodes and count as run.

The first protocol reruns the serving mug cell on the same eight seeds (67--101) with three references
(Table~\ref{tabC:follow}a). (J)~The joint-space person-blind control enters the 0.10\,m band on 2/64 episodes, within
the predicted bound of 6, and each policy is higher ($+40$, $+50$ and $+71$ points; Holm $p=0.0005$ for each). Unlike
the T1 excess, the T2 excess does not depend on the Cartesian control. This control delivers on only 7/64 episodes.
(H)~With the bystander not rendered but still scored, \pifive{} and \pifast{} do not change detectably, as predicted.
\grd{} does: 43/58 with the bystander rendered against 24/58 hidden ($+33$ points, Holm $p=0.006$). It enters the band
more often with the bystander in view, not less. (M)~With the capsule replaced by a human mesh no policy changes
detectably ($-13$, $+3$ and $+10$ points; Holm $p\geq 0.57$). Arms of eight cells detect differences of 22 to 32
points with 80\,\% power, so smaller effects of the bystander's appearance are not excluded.

The second protocol asks whether the tilt raised by ``Do not spill the coffee.'' (Sec.~\ref{sec:findings}) depends on
how the person is rendered or on pinned pick and place spots: \pifive{} at the dining table, person on the right, with
and without the clause, capsule or human mesh, pinned or default spots, four fresh seeds (103--113). The clause raises
the T4 rate in all four conditions (Table~\ref{tabC:follow}b): 90/116 against 3/120 pooled ($+75$ points; stratified
cell permutation, one-sided $p=0.00005$). Neither moderator is detectable (mesh minus capsule $-3$ points, default
minus pinned $-10$; Holm $p\geq 0.40$), and all five predictions held. With the clause the payload is delivered on
39/122 attempts, against 117/123 without it.

\begin{table}[!htb]
\centering
\caption{\textbf{The serving-cell excess holds against a joint-space control and with a human mesh; the spill-clause
effect holds with a mesh and unpinned spots.} (a)~Episodes with a link within 0.10\,m of the bystander's body\,/\,all
episodes at the pre-registered serving cell, eight cells per arm. Rendered: the capsule arms of
Table~\ref{tabC:prereg} (\grd{}: 35\,s). Hidden: not rendered, still scored. Last columns: policy minus joint-space
control with the capsule rendered, points; exact cell permutation $p$, Holm over the three. (b)~Carries tilted past
$45^{\circ}$\,/\,scored carries, \pifive{}, four cells per entry; difference in points.}
\label{tabC:follow}
\scriptsize
\setlength{\tabcolsep}{3pt}
\begin{tabular}{@{}l c c c r l@{}}
\toprule
\multicolumn{6}{@{}l}{\textit{(a) Serving cell, T2}}\\
Arm & Capsule rendered & Hidden & Human mesh & vs joint-space & $p$ (Holm) \\
\midrule
\pifive & 26/60 & 22/64 & 36/64 & $+40$ & 0.0005 (0.0005) \\
\pifast & 34/64 & 40/59 & 32/64 & $+50$ & 0.0002 (0.0005) \\
\grd & 43/58 & 24/58 & 41/64 & $+71$ & 0.0002 (0.0005) \\
Cartesian control & 0/64 & & & & \\
Joint-space control & 2/64 & & & & \\
\midrule
\multicolumn{6}{@{}l}{\textit{(b) Spill clause, T4}}\\
Condition & With the clause & Without & & Difference & \\
\midrule
Capsule, pinned spots & 27/32 & 1/27 & & $+81$ & \\
Human mesh, pinned spots & 24/30 & 1/32 & & $+77$ & \\
Capsule, default spots & 22/30 & 0/31 & & $+73$ & \\
Human mesh, default spots & 17/24 & 1/30 & & $+68$ & \\
\bottomrule
\end{tabular}
\end{table}
